% =============================================================================
% ПРИЛОЖЕНИЕ B: ПОЛНАЯ ФОРМАЛЬНАЯ МОДЕЛЬ ПАРАДОКСА ВРЕМЕННОЙ КООРДИНАЦИИ
% =============================================================================

\documentclass[12pt]{article}
\usepackage[a4paper,left=2cm,right=2cm,top=2.5cm,bottom=2.5cm]{geometry}
\usepackage{amsmath, amssymb, amsthm}
\usepackage{hyperref}
\usepackage{titlesec}
\usepackage{graphicx}
\usepackage{tikz}
\usepackage{caption}
\usepackage{booktabs}
\usepackage{array}
\usepackage{fontspec}
% Подключаем русский язык
\usepackage{polyglossia}
\setmainlanguage{russian}
\setotherlanguage{english}
\setmainfont{Times New Roman}
\newfontfamily{\cyrillicfont}{Times New Roman}
\newfontfamily{\cyrillicfontsf}{Arial}
\newfontfamily{\cyrillicfonttt}{Courier New}

% Все необходимые библиотеки TikZ
\usetikzlibrary{positioning, arrows.meta, shapes.geometric, calc}

\titleformat{\section}{\large\bfseries}{\thesection}{1em}{}
\titleformat{\subsection}{\normalsize\bfseries}{\thesubsection}{1em}{}
\titleformat{\subsubsection}{\normalsize\bfseries}{\thesubsubsection}{1em}{}

\renewcommand{\baselinestretch}{1.1}

\newtheorem{definition}{Определение}
\newtheorem{theorem}{Теорема}
\newtheorem{lemma}{Лемма}
\newtheorem{corollary}{Следствие}
\newtheorem{axiom}{Аксиома}

\begin{document}
	\begin{titlepage}
		\begin{center}
			\vspace*{0cm}
			\Huge\textbf{Приложение B. YUCT\\
				ПАРАДОКС ВРЕМЕННОЙ КООРДИНАЦИИ}
			
			\vspace{0cm}
			\LARGE\textit{От информационно-теоретических основ к экспериментальной верификации \\
				Математическое разрешение парадокса $K_{\text{eff}} > 1$}
			
			\vspace{1cm}
			\Large
			Алексей В. Якушев\\
			Yakushev Research\\
			Теория YUCT: \url{https://yuct.org/}\\
			Протокол YPSDC: \url{https://ypsdc.com/}\\
			
			\vspace{2cm}
			\textbf DOI: \url{https://doi.org/10.5281/zenodo.18444598}\\
			
			\vspace{1cm}
			
			\vspace{0cm}
			\large 
			Краткая версия этой работы была ранее опубликована и доступна по адресу\\ \url{https://doi.org/10.5281/zenodo.18362308}\\
			\vspace{1cm}
			Март 2026 \\ \large V.2.1.
			
			\vspace{5cm}
			\textcopyright~2026 Yakushev Research. Все права защищены.
			
			\newpage
			\vfill
			\normalsize
			\begin{minipage}{0.8\textwidth}
				\centering
				\textbf{Аннотация:} \\ В этом документе представлена полная формальная математическая модель парадокса временной координации в Объединённой теории координации Якушева (YUCT). Парадокс возникает, когда эффективность координации $K_{\text{eff}}$ превышает 1, что, казалось бы, нарушает классические пределы передачи информации. Мы даём строгие определения, теоремы и доказательства, показывающие, как опережающее знание с помощью предварительных словарей обеспечивает координацию, которая кажется временным парадоксом, но остаётся согласованной с физическими законами. Основные результаты включают: (1) формальную теорему разделения пропускной способности канала и координационной способности, (2) математический вывод $K_{\text{eff}} = R \cdot \eta$, где $R = n/m$ — коэффициент сжатия знаний, (3) точные условия, при которых $K_{\text{eff}}$ может превышать 1 без нарушения причинности, (4) теоремы многоузловой координации для масштабируемых систем, (5) энтропийную интерпретацию эффективности координации и (6) экспериментальные протоколы для измерения $K_{\text{eff}}$ в реальных системах.
			\end{minipage}
			
			\vspace{0.5cm}
			\normalsize
			\textbf{Ключевые слова:} парадокс временной координации, формальная модель YUCT, $K_{\text{eff}} > 1$, опережающее знание, предварительные словари, информационно-теоретическая координация, разделение пропускной способности канала, протокол YPSDC, сохранение причинности, энтропия координации.
			
			\vspace{0cm}
			\normalsize
			
			\textbf{Область применения:} формальные математические основы временной координации \\
			\textbf{Ключевые теоремы:} 8 формальных теорем с полными доказательствами \\
			\textbf{Экспериментальные протоколы:} 3 метода измерения $K_{\text{eff}}$ \\
			
			\vspace{0.5cm}
			\normalsize
			
		\end{center}
	\end{titlepage}	
	\begin{center}
		\vspace*{0.5cm}
		\Huge\textbf{Парадокс временной координации}
		\vspace{1cm}
	\end{center}
	
	\section{Полная формальная модель парадокса временной координации}
	\label{app:formal-model}
	
	\subsection{Формальные определения}
	\label{app:definitions}
	
	\subsubsection{Базовые объекты}
	
	Определим систему координации как кортеж $\mathcal{S} = (\mathcal{A}, \mathcal{O}, \mathcal{D}, \mathcal{C}, \mathcal{T})$, где:
	
	\begin{itemize}
		\item $\mathcal{A} = \{a_1, a_2, \dots, a_N\}$ — конечное множество возможных действий (активаций знаний),
		\item $\mathcal{O} = \{O_1, O_2, \dots, O_M\}$ — множество наблюдателей (узлов),
		\item $\mathcal{D} = \{\kappa_i \to a_i\}_{i=1}^N$ — предварительный словарь (биективное отображение кодов в действия),
		\item $\mathcal{C}$ — физический канал передачи с заданными параметрами,
		\item $\mathcal{T}$ — набор временных ограничений.
	\end{itemize}
	
	\subsubsection{Информационно-теоретические меры}
	
	\begin{definition}[Информационное содержание действия]
		Для каждого действия $a \in \mathcal{A}$ его полное описание требует:
		\[
		I(a) = n \ \text{бит},
		\]
		где $n$ — минимальное количество бит, необходимое для однозначного описания $a$ без предварительных знаний.
	\end{definition}
	
	\begin{definition}[Длина кода]
		Для каждого кода $\kappa \in \mathcal{K}$ (множества кодов в $\mathcal{D}$):
		\[
		|\kappa| = m \ \text{бит}, \quad \text{причём} \ m \ll n.
		\]
	\end{definition}
	
	\begin{definition}[Коэффициент сжатия знаний]
		Коэффициент сжатия знаний определяется как:
		\[
		R = \frac{n}{m} \gg 1.
		\]
	\end{definition}
	
	\subsection{Модель физического канала}
	\label{app:channel-model}
	
	\subsubsection{Фундаментальные ограничения}
	
	\begin{axiom}[Предел скорости света]
		Для любых двух узлов $O_i, O_j \in \mathcal{O}$, разделённых расстоянием $L_{ij}$, скорость распространения сигнала удовлетворяет:
		\[
		v_{\text{signal}} \leq c,
		\]
		где $c$ — скорость света в вакууме.
	\end{axiom}
	
	\begin{axiom}[Пропускная способность канала]
		Канал $\mathcal{C}$ имеет максимальную пропускную способность:
		\[
		C_{\text{max}} = \lim_{T \to \infty} \frac{1}{T} \max_{X} I(X; Y) \quad \text{[бит/с]},
		\]
		где $X$ — входной сигнал, $Y$ — выходной сигнал.
	\end{axiom}
	
	\subsubsection{Модель времени передачи}
	
	Время, необходимое для передачи сообщения из $I$ бит:
	\[
	T_{\text{transmit}}(I) = \frac{I}{C_{\text{eff}}} + \frac{L}{c} + \tau_{\text{proc}},
	\]
	где $C_{\text{eff}} \leq C_{\text{max}}$ — эффективная пропускная способность канала, а $\tau_{\text{proc}}$ — задержка обработки.
	
	\subsection{Процессы координации}
	\label{app:coordination-processes}
	
	\subsubsection{Наивная координация (без словаря)}
	
	\textbf{Сценарий 1:} Узел $O_1$ хочет, чтобы узел $O_2$ выполнил действие $a$.
	
	\begin{enumerate}
		\item $O_1$ формулирует полное описание $a$: $I(a) = n$ бит.
		\item $O_1$ передаёт это описание $O_2$: время $T_1 = T_{\text{transmit}}(n)$.
		\item $O_2$ декодирует и понимает действие: время $T_2 = \alpha n$, где $\alpha > 0$.
		\item $O_2$ выполняет действие: время $T_3$.
		\item $O_2$ отправляет подтверждение: время $T_4 = T_{\text{transmit}}(p)$, где $p$ — размер подтверждения.
	\end{enumerate}
	
	Общее наивное время координации:
	\[
	T_{\text{naive}}(a) = T_1 + T_2 + T_3 + T_4.
	\]
	
	\begin{theorem}[Нижняя граница для наивной координации]
		\label{thm:naive-lower-bound}
		Для любой пары узлов $O_1, O_2$ и действия $a$ наивное время координации ограничено снизу:
		\[
		T_{\text{naive}}(a) \geq \frac{n + p}{C_{\text{eff}}} + \frac{2L}{c} + \alpha n.
		\]
	\end{theorem}
	
	\subsubsection{Координация на основе словаря (YPSDC)}
	
	\textbf{Сценарий 2:} Узлы $O_1$ и $O_2$ совместно используют предварительный словарь $\mathcal{D}$.
	
	\begin{enumerate}
		\item $O_1$ выбирает код $\kappa$, соответствующий действию $a$: $|\kappa| = m$ бит.
		\item $O_1$ передаёт $\kappa$ узлу $O_2$: время $T_1' = T_{\text{transmit}}(m)$.
		\item $O_2$ находит действие в $\mathcal{D}$: время $T_2' = \beta m$, причём $\beta \ll \alpha$.
		\item $O_2$ выполняет действие: время $T_3' = T_3$.
	\end{enumerate}
	
	Критическое наблюдение: $O_1$ может действовать так, \emph{как если бы} $O_2$ уже выполнил $a$, сразу после отправки $\kappa$.
	
	Эффективное время координации со словарём:
	\[
	T_{\text{coord}}(a) = T_1' + T_2'.
	\]
	
	\begin{figure}[htbp]
		\centering
		\begin{tikzpicture}[node distance=2cm, auto, >=Stealth]
			\node (O1) [draw, rectangle, rounded corners, minimum width=2cm, minimum height=1cm] {Наблюдатель 1};
			\node (O2) [draw, rectangle, rounded corners, minimum width=2cm, minimum height=1cm, right=of O1] {Наблюдатель 2};
			\node (dict) [draw, ellipse, below=2cm of $(O1)!0.5!(O2)$] {Предварительный словарь $\mathcal{D}$};
			\draw[->, thick] (O1) -- node[above] {короткий код $\kappa$} (O2);
			\draw[->, dashed] (O1) -- node[left, near start] {совместное} (dict);
			\draw[->, dashed] (O2) -- node[right, near start] {использование} (dict);
			\node[below=0.3cm of O1, align=center] {в $t_0$:\\ знает, что $a$ будет выполнено};
			\node[below=0.3cm of O2, align=center] {в $t_0 + L/c$:\\ выполняет $a$};
		\end{tikzpicture}
		\caption{Базовая схема координации YPSDC. Опережающее знание возникает в момент $t_0$ из общего словаря.}
		\label{fig:basic}
	\end{figure}
	
	\subsection{Парадокс временной координации}
	\label{app:paradox}
	
	\begin{definition}[Опережающее знание]
		В момент $t_0$ отправки кода $\kappa$ узел $O_1$ обладает опережающим знанием о будущем действии узла $O_2$:
		\[
		K_{\text{advance}}(O_1, O_2, a, t_0) = \text{Истина}
		\]
		тогда и только тогда, когда существует общий словарь $\mathcal{D}$ такой, что $\mathcal{D}(\kappa) = a$.
	\end{definition}
	
	\begin{lemma}[Существование опережающего знания]
		\label{lem:advance-knowledge}
		Для любых $O_1, O_2, a$ с общим словарём $\mathcal{D}$:
		\[
		K_{\text{advance}}(O_1, O_2, a, t_0) = \text{Истина}
		\]
		в момент $t_0$, когда отправляется $\kappa$.
	\end{lemma}
	
	\begin{proof}
		По построению $\mathcal{D}$, отправка $\kappa$ гарантирует, что $O_2$ выполнит $a$ в момент времени $t_0 + T_{\text{coord}}(a)$. Следовательно, в момент $t_0$ $O_1$ может с уверенностью предсказать будущее состояние $O_2$.
	\end{proof}
	
	\subsection{Информационные ёмкости координации}
	\label{app:capacities}
	
	\begin{definition}[Информационная ёмкость канала]
		\[
		C_{\text{channel}} = C_{\text{eff}} \quad \text{[бит/с]}.
		\]
	\end{definition}
	
	\begin{definition}[Координационная информационная ёмкость]
		\[
		C_{\text{coord}} = \lim_{T \to \infty} \frac{1}{T} \sum_{t=1}^{T} I(a_t) \quad \text{[бит/с]},
		\]
		где $a_t$ — действие, активированное в момент времени $t$.
	\end{definition}
	
	\begin{theorem}[Теорема разделения ёмкостей]
		\label{thm:capacity-separation}
		Для системы координации $\mathcal{S}$ с коэффициентом сжатия знаний $R = n/m$:
		\[
		C_{\text{coord}} = R \cdot C_{\text{channel}} \cdot \eta,
		\]
		где $\eta = \frac{T_{\text{channel}}}{T_{\text{coord}}} \leq 1$ — фактор эффективности времени.
	\end{theorem}
	
	\begin{proof}
		В течение интервала времени $\Delta T$ канал может передать:
		\[
		N_{\text{codes}} = \frac{C_{\text{channel}} \cdot \Delta T}{m} \quad \text{кодов}.
		\]
		Каждый код активирует действие с информационным содержанием $n$ бит, так что общая координированная информация составляет:
		\[
		I_{\text{total}} = N_{\text{codes}} \cdot n = \frac{C_{\text{channel}} \cdot \Delta T}{m} \cdot n.
		\]
		Следовательно,
		\[
		C_{\text{coord}} = \frac{I_{\text{total}}}{\Delta T} = C_{\text{channel}} \cdot \frac{n}{m} = R \cdot C_{\text{channel}}.
		\]
		Учёт временных задержек (например, обработка, распространение) вводит фактор эффективности $\eta$.
	\end{proof}
	
	\begin{figure}[htbp]
		\centering
		\begin{tikzpicture}[node distance=2.5cm, auto, >=Stealth]
			\node (O1) [draw, rectangle, rounded corners, minimum width=2.2cm, minimum height=1.2cm] {Наблюдатель 1 $O_1$};
			\node (O2) [draw, rectangle, rounded corners, minimum width=2.2cm, minimum height=1.2cm, right=of O1] {Наблюдатель 2 $O_2$};
			\node (dict) [draw, ellipse, below=2.5cm of $(O1)!0.5!(O2)$] {Предварительный словарь $\mathcal{D} = \{\kappa_i \to \mathcal{A}_i\}$};
			\draw[->, thick] (O1) -- node[above] {$\kappa_1$} (O2);
			\draw[->, thick] (O2) -- node[below] {$\kappa_2$ (подтверждение)} (O1);
			\draw[->, dashed] (O1) -- (dict);
			\draw[->, dashed] (O2) -- (dict);
			\node[below=0.2cm of O1, align=center] {Знает $\mathcal{A}_2$ в $t_0$};
			\node[below=0.2cm of O2, align=center] {Выполняет $\mathcal{A}_2$ в $t_0+L/c$};
		\end{tikzpicture}
		\caption{Дуплексная координация с предварительным словарём. Оба узла совместно используют один словарь, что обеспечивает двунаправленное опережающее знание.}
		\label{fig:duplex}
	\end{figure}
	
	\subsection{Фактор эффективности координации \(K_{\mathrm{eff}}\) и разделение ёмкостей}
	
	\begin{definition}[Эффективность координации]
		\emph{Эффективность координации} YPSDC-системы равна
		\begin{equation}
			K_{\mathrm{eff}} = \frac{H(\mathcal{A})}{H(\mathcal{K})},
			\label{eq:Kdef}
		\end{equation}
		где \(H(\mathcal{A})\) — энтропия пространства действий, а \(H(\mathcal{K})\) — энтропия распределения индексов.
	\end{definition}
	
	\begin{definition}[Координационная ёмкость]
		\emph{Координационная ёмкость} \(C_{\mathrm{coord}}\) — это максимальная скорость, с которой получатель может извлекать осмысленную информацию из словаря:
		\begin{equation}
			C_{\mathrm{coord}} = \lim_{\Delta t \to \infty} \frac{1}{\Delta t} \sum_{t=1}^{N_{\mathrm{actions}}} I(A_{i_t}),
		\end{equation}
		где \(A_{i_t}\) — действие, активированное во время $t$-й передачи.
	\end{definition}
	
	\begin{theorem}[Разделение ёмкостей (вторая форма)]
		\label{thm:capacity-separation}
		Для YPSDC-системы с эффективностью координации \(K_{\mathrm{eff}}\) и фактором накладных расходов \(\eta \le 1\) координационная ёмкость и пропускная способность канала связаны соотношением
		\begin{equation}
			\boxed{C_{\mathrm{coord}} = K_{\mathrm{eff}} \cdot C_{\mathrm{channel}} \cdot \eta}.
			\label{eq:capacity-separation}
		\end{equation}
	\end{theorem}
	
	\begin{proof}
		За интервал времени \(\Delta T\) канал может передать не более \(C_{\mathrm{channel}} \Delta T\) бит. Для индексов фиксированной длины \(\ell\) число переданных индексов равно \(N = \lfloor C_{\mathrm{channel}} \Delta T / \ell \rfloor\). Каждый индекс активирует действие, несущее в среднем \(H(\mathcal{A})\) бит информации. Общая осмысленная информация составляет \(I_{\mathrm{total}} = N \cdot H(\mathcal{A}) \approx \frac{C_{\mathrm{channel}} \Delta T}{\ell} H(\mathcal{A})\). Средняя скорость равна \(C_{\mathrm{coord}} = \frac{I_{\mathrm{total}}}{\Delta T} = C_{\mathrm{channel}} \cdot \frac{H(\mathcal{A})}{\ell} = C_{\mathrm{channel}} \cdot K_{\mathrm{eff}}\). Фактор \(\eta\) учитывает накладные расходы на обработку и доступ к словарю.
	\end{proof}
	
	\begin{theorem}[Обратная теорема]
		Для любой схемы, использующей фиксированный словарь \(D\) размера \(M = 2^\ell\) и передающей индексы по DMC с пропускной способностью \(C_{\mathrm{channel}}\), координационная ёмкость удовлетворяет
		\[
		C_{\mathrm{coord}} \le K_{\mathrm{eff}} \cdot C_{\mathrm{channel}}.
		\]
	\end{theorem}
	
	\begin{proof}
		Пусть источник индексов имеет энтропию \(H(\mathcal{K})\). По неравенству обработки данных, \(I(\kappa;\hat{\kappa}) \le I(X;Y) \le C_{\mathrm{channel}}\). Для надёжной работы индекс должен декодироваться с исчезающей ошибкой, что требует скорости индексов \(R_{\mathcal{K}} \le C_{\mathrm{channel}}\) (теорема Шеннона о кодировании для шумного канала). Средняя скорость осмысленной информации равна \(R_{\mathcal{A}} = R_{\mathcal{K}} \cdot K_{\mathrm{eff}} \le K_{\mathrm{eff}} \cdot C_{\mathrm{channel}}\). Следовательно, \(C_{\mathrm{coord}} \le K_{\mathrm{eff}} \cdot C_{\mathrm{channel}}\).
	\end{proof}
	
	\subsection{Энтропийная интерпретация координации}
	\label{app:entropy}
	
	\begin{definition}[Энтропия словаря]
		Энтропия словаря $\mathcal{D}$:
		\[
		H(\mathcal{D}) = - \sum_{i=1}^{N} p_i \log_2 p_i,
		\]
		где $p_i$ — вероятность использования действия $a_i$.
	\end{definition}
	
	\begin{lemma}[Максимальная координационная ёмкость]
		Максимальная координационная ёмкость ограничена:
		\[
		C_{\text{coord}}^{\text{max}} = H(\mathcal{D}) \cdot \nu_{\text{max}},
		\]
		где $\nu_{\text{max}} = 1 / T_{\text{coord}}^{\text{min}}$ — максимальная частота координации.
	\end{lemma}
	
	\begin{definition}[Энтропия координации]
		Изменение энтропии системы вследствие координации:
		\[
		\Delta S_{\text{coord}} = k_B \ln(K_{\mathrm{eff}}),
		\]
		где $k_B$ — постоянная Больцмана.
	\end{definition}
	
	Интерпретация: Увеличение $K_{\mathrm{eff}}$ соответствует уменьшению энтропии (увеличению порядка) координированной системы.
	
	\subsection{Проверяемые предсказания формальной модели}
	\label{app:predictions}
	
	\begin{enumerate}
		\item \textbf{Квантование $K_{\mathrm{eff}}$:} Для оптимально спроектированных систем $K_{\mathrm{eff}} = 2^k$, где $k \in \mathbb{N}$, и в идеальных условиях $k = \log_2 R$.
		
		\item \textbf{Масштабирование с размером системы:} Для системы из $M$ узлов $K_{\mathrm{eff}}(M) \propto M \log_2 R$.
		
		\item \textbf{Фундаментальный предел:} Существует фундаментальный предел:
		\[
		K_{\mathrm{eff}}^{\text{max}} = \frac{c}{\Delta x \cdot \Delta t},
		\]
		где $\Delta x$ — минимальное пространственное разрешение системы, а $\Delta t$ — минимальное временное разрешение.
	\end{enumerate}
	
	\subsection{Экспериментальный протокол измерения $K_{\mathrm{eff}}$}
	\label{app:experiment}
	
	\begin{enumerate}
		\item \textbf{Измерение $C_{\text{channel}}$:} Используйте стандартные инструменты измерения сети (например, iperf, ping) для оценки эффективной пропускной способности канала.
		
		\item \textbf{Измерение $C_{\text{coord}}$:}
		\begin{itemize}
			\item Определите репрезентативный набор действий $\mathcal{A}$.
			\item Измерьте общее время $T_{\text{total}}$ для выполнения серии координированных действий.
			\item Вычислите $C_{\text{coord}} = \frac{\sum I(a_i)}{T_{\text{total}}}$.
		\end{itemize}
		
		\item \textbf{Вычисление $K_{\mathrm{eff}}$:}
		\[
		K_{\mathrm{eff}} = \frac{C_{\text{coord}}}{C_{\text{channel}}}.
		\]
	\end{enumerate}
	
	Ожидаемые диапазоны для различных систем:
	\begin{itemize}
		\item Системы команд человека: $K_{\mathrm{eff}} \approx 10^2 - 10^3$.
		\item Компьютерные сети: $K_{\mathrm{eff}} \approx 10^3 - 10^6$.
		\item Биологические системы: $K_{\mathrm{eff}} \approx 10^6 - 10^9$.
	\end{itemize}
	
	\begin{table}[htbp]
		\centering
		\caption{Фундаментальное разделение: передача данных vs. координация состояний.}
		\label{tab:separation}
		\begin{tabular}{lp{6cm}p{6cm}}
			\toprule
			\textbf{Аспект} & \textbf{Передача данных} & \textbf{Координация} \\
			\midrule
			Физический сигнал & Да & Нет (знание состояния) \\
			Предел скорости & $v \leq c$ & $v_{\text{coord}} \to \infty$* \\
			Информация & $I > 0$ & $K_{\text{eff}} > I$ \\
			Энергия & $E > 0$ & $\Delta S_{\text{coord}}$ \\
			\bottomrule
			\multicolumn{3}{l}{*в смысле мгновенного согласования состояний через предварительные знания.}
		\end{tabular}
	\end{table}
	
	\subsection{Заключение формальной модели}
	
	Представленная в этом приложении формальная математическая модель строго устанавливает:
	
	\begin{enumerate}
		\item Информационная ёмкость координации $C_{\text{coord}}$ и пропускная способность канала $C_{\text{channel}}$ — различные физические величины.
		
		\item Парадокс времени координации: априорные словари позволяют узлу иметь опережающее знание о будущих действиях других узлов до того, как они будут фактически выполнены.
		
		\item Количественное соотношение:
		\[
		K_{\mathrm{eff}} = \frac{C_{\text{coord}}}{C_{\text{channel}}} = R \cdot \eta \gg 1,
		\]
		где $R = n/m \gg 1$ — коэффициент сжатия знаний.
		
		\item Фундаментальное ограничение: рост $K_{\mathrm{eff}}$ ограничен только сложностью создания и поддержания словаря, а не физическими законами передачи информации.
	\end{enumerate}
	
	Эта модель обеспечивает строгую математическую основу для анализа и проектирования высокоэффективных систем координации в физике, биологии, социологии и технологиях.
	
	
	\section{Экспериментальные протоколы и мысленный эксперимент для парадокса временной координации}
	\label{app:experimental-protocols}
	
	Формальная модель парадокса временной координации (разделы \ref{app:formal-model}--\ref{app:keff}) обеспечивает ясную математическую основу, но её практическая проверка требует конкретных экспериментальных протоколов. В этом разделе мы предлагаем несколько экспериментов — от простого мысленного (Gedanken) эксперимента до реализуемых лабораторных тестов — которые напрямую измеряют эффективность координации \(K_{\mathrm{eff}}\) и демонстрируют парадокс опережающего знания.
	
	\subsection{Мысленный эксперимент: аналогия с 2FA (двухфакторной аутентификацией)}
	\label{subsec:thought-2fa}
	
	Рассмотрим двух наблюдателей, Алису и Боба, которые совместно используют криптографический словарь \( \mathcal{D} \) (например, одноразовый блокнот из \(M\) предварительно разделённых ключей). Алиса хочет, чтобы Боб выполнил одно из \(M\) возможных действий (например, перевёл сумму денег, запустил ракету или отправил конкретное сообщение). В наивном протоколе Алиса отправила бы полное описание действия, требующее \(n\) бит. С общим словарём она отправляет только короткий индекс \(\kappa\) длиной \(\ell = \log_2 M\) бит.
	
	\textbf{Установка:}
	\begin{itemize}
		\item Алиса и Боб находятся на известном расстоянии \(L\) (например, 1 км).
		\item Канал связи имеет измеренную полосу пропускания и задержку распространения \(L/c\).
		\item Словарь \(\mathcal{D}\) распространяется офлайн через доверенного курьера или по предварительному защищённому каналу (занимает произвольно долгое время, но завершается до эксперимента).
		\item Третий наблюдатель, Чарли, оснащённый двумя синхронизированными атомными часами, записывает точное время, когда Алиса инициирует передачу, и когда Боб завершает действие.
	\end{itemize}
	
	\textbf{Процедура:}
	\begin{enumerate}
		\item Алиса выбирает действие \(a_i\) с полной длиной описания \(n\).
		\item В момент времени \(t_0\) она отправляет соответствующий индекс \(\kappa_i\) Бобу.
		\item Индекс распространяется со скоростью \(c\) и прибывает к Бобу в момент времени \(t_0 + L/c + \tau_{\text{tx}}\), где \(\tau_{\text{tx}} = \ell / C_{\text{eff}}\).
		\item Боб мгновенно находит действие в словаре (время поиска \(\tau_{\text{lookup}} \ll \tau_{\text{tx}}\)) и выполняет его в момент времени \(t_0 + L/c + \tau_{\text{tx}} + \tau_{\text{lookup}}\).
	\end{enumerate}
	
	\textbf{Наблюдение парадокса:}
	В момент \(t_0\) Алиса \emph{знает}, что Боб выполнит действие \(a_i\) в будущий момент времени \(t_0 + L/c + \tau_{\text{tx}} + \tau_{\text{lookup}}\). Это опережающее знание основано не на сверхсветовой сигнализации, а на предварительно разделённом словаре. Чарли, у которого нет словаря, не может предсказать действие до тех пор, пока индекс не будет передан. Таким образом, парадокс разрешается разделением информации о словаре (офлайн) и индекса (онлайн).
	
	\textbf{Количественная мера:}
	Эффективность координации равна
	\[
	K_{\mathrm{eff}} = \frac{n}{\ell} \cdot \frac{\tau_{\text{naive}}}{\tau_{\text{actual}}},
	\]
	где \(\tau_{\text{naive}}\) — время, необходимое для отправки полного описания (включая распространение и обработку). В пределе быстрого поиска и пренебрежимо малого времени передачи индекса \(K_{\mathrm{eff}} \approx n/\ell\), что может быть сделано сколь угодно большим за счёт увеличения размера словаря.
	
	\subsection{Реализуемый лабораторный эксперимент: связь с отложенным выбором}
	\label{subsec:lab-delayed-choice}
	
	\textbf{Цель:} Напрямую измерить \(K_{\mathrm{eff}}\) в управляемой цифровой системе связи и подтвердить, что \(K_{\mathrm{eff}} > 1\) может быть достигнуто без нарушения причинности.
	
	\textbf{Аппаратура:}
	\begin{itemize}
		\item Два компьютера (Алиса и Боб), соединённые Ethernet-линией с известной пропускной способностью и переменной линией задержки (например, длинное оптическое волокно) для имитации расстояния.
		\item Предварительно разделённый словарь из \(M\) сообщений, каждое длиной \(n\) бит (например, \(M = 2^{20}\), \(n = 10^6\) бит).
		\item Высокоточные метки времени (наносекундное разрешение) с использованием GPS-синхронизированных часов.
		\item Программное обеспечение для измерения времени между отправкой индекса и получением подтверждения о выполненном действии.
	\end{itemize}
	
	\textbf{Протокол:}
	\begin{enumerate}
		\item Этап калибровки: измерьте время полного цикла для отправки полного сообщения (наивный протокол) по той же линии. Запишите \(T_{\text{naive}}\).
		\item Этап координации: для каждого испытания Алиса случайным образом выбирает сообщение, отправляет его индекс и записывает время \(T_{\text{coord}}\) до тех пор, пока Боб не подтвердит выполнение.
		\item Повторите для различных размеров словаря \(M\) и длин сообщений \(n\). Вычислите \(K_{\mathrm{eff}} = T_{\text{naive}} / T_{\text{coord}}\).
		\item Также измерьте пропускную способность канала \(C_{\text{channel}}\) и время поиска в словаре, чтобы проверить теорему разделения ёмкостей.
	\end{enumerate}
	
	\textbf{Ожидаемый результат:} Для достаточно больших словарей \(K_{\mathrm{eff}}\) значительно превысит 1, демонстрируя, что эффективная скорость передачи осмысленной информации может превышать пропускную способность канала. Эксперимент также подтверждает, что распространение индекса подчиняется \(v \le c\); причинность сохраняется, поскольку словарь был распространён офлайн.
	
	\subsection{Эксперимент по социальной координации: время реакции человека с общим контекстом}
	\label{subsec:social-experiment}
	
	\textbf{Цель:} Продемонстрировать парадокс временной координации в социальной среде человека, по аналогии с мысленным экспериментом 2FA.
	
	\textbf{Установка:}
	\begin{itemize}
		\item Две группы испытуемых (например, студентов) обучаются набору из \(M\) сложных команд (например, «нарисовать дом», «решить головоломку», «прочитать стихотворение»). Каждая команда требует нескольких секунд на выполнение, а её полное описание (текст + инструкции) составляет около \(n\) бит.
		\item Этап обучения создаёт общий «словарь».
		\item Во время эксперимента одна группа («отправитель») видит команду и должна передать её другой группе («получатель»), используя только одно короткое слово (индекс), о котором они договорились во время обучения.
		\item Расстояние между группами варьируется (например, разные комнаты, здания), чтобы ввести измеримую задержку распространения звуковых или световых сигналов.
		\item Измеряется время от момента, когда отправитель получает команду, до момента, когда получатель завершает действие.
	\end{itemize}
	
	\textbf{Контроль:} Те же самые команды передаются наивным способом (например, путём подробного описания полного действия) без предварительного обучения.
	
	\textbf{Измерения:} Сравните время координации с предварительным словарём и без него. Вычислите \(K_{\mathrm{eff}}\) как отношение наивного времени ко времени координации. Покажите, что даже если индекс распространяется с конечной скоростью (скорость звука или света), эффективная координация может выглядеть почти мгновенной по сравнению с наивным методом.
	
	\textbf{Ожидаемый результат:} \(K_{\mathrm{eff}}\) будет масштабироваться со сложностью команд и размером общего словаря, подтверждая теоретические предсказания.
	
	\subsection{Экспериментальная проверка теоремы разделения ёмкостей}
	\label{subsec:capacity-theorem}
	
	Теорема разделения ёмкостей (Теорема \ref{thm:capacity-separation}) утверждает, что \(C_{\text{coord}} = K_{\mathrm{eff}} \cdot C_{\text{channel}} \cdot \eta\). Для её проверки необходимо независимо измерить как \(C_{\text{coord}}\), так и \(C_{\text{channel}}\).
	
	\textbf{Метод:}
	\begin{itemize}
		\item Используйте ту же установку, что и в разделе \ref{subsec:lab-delayed-choice}, но меняйте полосу пропускания канала (например, с помощью искусственного ограничения скорости).
		\item Для каждого значения полосы пропускания измерьте максимальную скорость, с которой индексы могут быть надёжно переданы (\(C_{\text{channel}}\)).
		\item Одновременно измерьте скорость, с которой завершаются осмысленные действия (\(C_{\text{coord}}\)), подсчитывая количество успешно активированных действий в секунду.
		\item Постройте график зависимости \(C_{\text{coord}}\) от \(C_{\text{channel}}\) для разных размеров словаря. Наклон должен быть равен \(K_{\mathrm{eff}}\), а пересечение с осью должно быть нулевым (с точностью до фактора эффективности времени \(\eta\)).
	\end{itemize}
	
	\subsection{Интерпретация и значение}
	
	Эти эксперименты напрямую демонстрируют основную идею YPSDC: разделение офлайн-распространения словаря и онлайн-передачи индекса позволяет достичь эффективности координации \(K_{\mathrm{eff}} > 1\) без нарушения причинности. Парадокс опережающего знания разрешается, потому что словарь разделён заранее, а сам индекс не несёт новой информации, кроме указания на существующую запись.
	
	Более того, эксперименты предоставляют количественные меры \(K_{\mathrm{eff}}\) в контролируемых условиях, что позволяет калибровать теорию и сравнивать её с реальными системами (например, интернет-протоколами, финансовыми сетями, биологической сигнализацией). Они также служат образовательным инструментом для иллюстрации контр-интуитивной, но каузальной природы высокоэффективной координации.
	
	\subsection{Заключение экспериментального раздела}
	
	Реализуя эти протоколы — от простого мысленного эксперимента до лабораторных и социальных тестов — исследователи могут эмпирически подтвердить парадокс временной координации и теорему разделения ёмкостей. Результаты подтвердят, что \(K_{\mathrm{eff}}\) — это не просто теоретическая конструкция, а измеримая величина, которая может превышать единицу, прокладывая путь для практических приложений в связи, распределённых вычислениях и коллективном принятии решений.
	
\end{document}